\documentclass{beamer}
\usetheme{default}

\title{Tides in Small Harbours}
\subtitle{What the gauges say}
\author{Ada Example \and Ben Sample}
\institute{Institute of Invented Places}
\date{\today}

\newcommand{\term}[1]{\alert{#1}}

\begin{document}

\begin{frame}
  \titlepage
\end{frame}

\begin{frame}{Outline}
  \tableofcontents[hideallsubsections]
\end{frame}

\section{Gauges}
\begin{frame}
  \frametitle{Three gauges}
  \framesubtitle{and what each one measures}
  \begin{itemize}
    \item<1-> A float in a well, read by hand.
    \item<2-> A pressure sensor on the sea bed.
    \item<3-| alert@3> A radar above the quay.
  \end{itemize}
  \pause
  Each has its own \term{drift}, and \only<3>{the radar's is the smallest}.
\end{frame}

\begin{frame}[fragile]{A reading}
  \begin{verbatim}
time  level
06:00 1.20
\end{verbatim}
\end{frame}

\section{Results}
\begin{frame}{Side by side}{two harbours}
  \begin{columns}[T]
    \begin{column}{0.48\textwidth}
      The north harbour fills in two hours:
      \[ h(t) = h_0 \sin(\omega t) \]
    \end{column}
    \begin{column}{0.48\textwidth}
      \begin{block}{Remember}
        The south one takes \textbf{three}.
      \end{block}
      \begin{alertblock}{Careful}
        Storms shift both.
      \end{alertblock}
      \begin{exampleblock}{An example}
        A spring tide in March.
      \end{exampleblock}
    \end{column}
  \end{columns}
\end{frame}

\begin{frame}{Steps}
  \begin{enumerate}[<+->]
    \item Read the gauge.
    \item Correct for the drift.
    \item Compare the harbours.
  \end{enumerate}
  \uncover<4->{And then publish.}
  \visible<4>{Only on the fourth slide.}
\end{frame}

\begin{frame}
  \centering
  \Large Thank you.\\[1em]
  \normalsize Questions?
\end{frame}

\end{document}
